% ═══════════════════════════════════════════════════════════════
%  appendices/c-glossary.tex
%  Приложение В. Глоссарий
% ═══════════════════════════════════════════════════════════════

\chapter{Глоссарий}
\label{app:glossary}

\begin{description}

\item[ADC] Analog-to-Digital Converter~--- аналого-цифровой
    преобразователь. Превращает напряжение в число.

\item[AppData] Центральная структура нашего приложения. Содержит
    всё состояние: указатели на виджеты, порт, буферы, настройки.

\item[AppImage] Формат упаковки приложений для Linux. Один файл,
    который содержит всё необходимое.

\item[Cairo] Векторная библиотека рисования. Используется в GTK
    для отрисовки графики. Работает через контекст \texttt{cairo\_t}.

\item[CMake] Система сборки. Описывает проект в файле
    \texttt{CMakeLists.txt}.

\item[DAC] Digital-to-Analog Converter~--- цифро-аналоговый
    преобразователь. Превращает число в напряжение.

\item[DDS] Direct Digital Synthesis~--- прямой цифровой синтез.
    Метод генерации сигналов через таблицу и аккумулятор фазы.

\item[DMA] Direct Memory Access~--- прямой доступ к памяти.
    Позволяет перекладывать данные без участия CPU.

\item[fd] File descriptor~--- файловый дескриптор. Целое число,
    которое ядро использует для ссылки на открытый файл или
    устройство.

\item[GIOChannel] Абстракция GLib для работы с потоками ввода-вывода.
    Позволяет встроить работу с портом в главный цикл GTK.

\item[GKeyFile] Парсер INI-файлов в составе GLib. Используется
    для хранения настроек.

\item[GLib] Базовая библиотека GNOME. Предоставляет структуры данных,
    таймеры, работу с файлами, парсеры.

\item[GMutex] Мьютекс в GLib. Защищает общие данные от
    одновременного доступа из разных потоков.

\item[GObject] Объектная модель на C. Основа GTK. Виджеты~---
    это GObject'ы.

\item[GTK] GIMP Toolkit~--- библиотека для создания графических
    интерфейсов.

\item[GThread] Поток в GLib. Позволяет выполнять код
    параллельно с главным потоком.

\item[gdb] GNU Debugger~--- пошаговый отладчик.

\item[HAL] Hardware Abstraction Layer~--- слой абстракции
    над железом. В нашем проекте~--- это \texttt{serial\_port.c}
    и \texttt{debug\_log.c}.

\item[libudev] Библиотека для работы с udev.

\item[LineBuffer] Накопитель строк. Собирает строки из потока
    байтов, разделённых \texttt{\textbackslash n} или
    \texttt{\textbackslash r}.

\item[OpenOCD] Open On-Chip Debugger~--- GDB-сервер для отладки
    STM32.

\item[Ping-pong] Двойная буферизация. Пока DMA пишет в одну
    половину буфера~--- CPU читает из другой.

\item[PlotData] Кольцевой буфер для хранения последних N
    измерений.

\item[POSIX] Portable Operating System Interface~--- стандарт
    системных вызовов в Unix-подобных ОС.

\item[PWM] Pulse-Width Modulation~--- широтно-импульсная
    модуляция. Быстрое переключение между 0 и 3.3 В.

\item[Sa/s] Samples per second~--- семплов в секунду. Единица
    измерения частоты дискретизации ADC.

\item[STM32] Микроконтроллеры на ядрах ARM Cortex-M от
    STMicroelectronics.

\item[termios] Структура и набор функций для настройки
    последовательного порта.

\item[udev] Менеджер устройств в Linux. Создаёт файлы в
    \texttt{/dev}, хранит свойства устройств, применяет правила.

\item[UART] Universal Asynchronous Receiver-Transmitter~---
    аппаратный модуль для последовательной связи.

\item[USB] Universal Serial Bus~--- универсальная
    последовательная шина.

\item[valgrind] Инструмент для поиска утечек памяти.

\item[VID:PID] Vendor ID и Product ID~--- идентификаторы
    USB-устройства. Уникальная пара для каждого типа устройств.

\item[VMIN/VTIME] Параметры \texttt{termios}, управляющие
    режимом чтения: сколько байт ждать и сколько времени.

\item[XDG] X Desktop Group~--- стандарт для хранения
    пользовательских данных. Определяет, где хранить настройки
    (\texttt{\$XDG\_CONFIG\_HOME}).

\end{description}
